\documentclass[12pt]{article}
%\renewcommand\normalsize{\fontsize{10pt}{12pt}\selectfont}
\usepackage{graphicx} % Required for inserting images
\usepackage[left=2cm,right=2cm,top=2cm,bottom=1.5cm]{geometry}

\title{PP}
%\author{Dennis Wilkinson}
\date{due Tuesday April 29}

\begin{document}

\setlength{\parindent}{0pt}
\setlength{\parskip}{3mm}

\underline{FS1 Wilkinson Final Exam June 12, 2025}

\textbf{1.} Two equal positive charges are fixed in space as below. X, Y, and Z are reference points.

\vspace{2cm}


\textbf{a.} Draw in the electric field lines around the charges. What direction is the field at point X?

\textbf{b.} \textbf{(i)} Which point has the \textit{lowest} voltage: X, Y, or Z? \textbf{(ii)} Would it take more energy to move a positive charge from far away to point X or to point Z? 

\textbf{c.} Say the charges are +8 nC, meaning $8 \times 10^{-9}$ C, they are 2 cm apart, and point Y is 2 cm left of the leftmost charge. A new +6 nC charge is placed at point Y. Find the force on it.

\textbf{d*.} The electric field of a single charge decreases as we move away from it, because of the $r^2$ in the denominator of $E=\frac{kQ}{r^2}$. In HW 6 we found that for a dipole, $E$ decreases more quickly (like $1/r^3$) as we move away. For this problem, how fast would $E$ would decrease as we move away - like for a single charge, like for a dipole, in between, something else? Defend your answer.

\bigskip

\textbf{2a.} Give two reasons why household circuits are wired in parallel (not series).

\textbf{b.} Pick \textbf{one} of the following and explain briefly how it works: speaker, microphone, cordless charger, induction cooktop. Include a diagram.

\bigskip

\textbf{3.} A rail gun consists of parallel rails, which are fixed in place, and a cross piece called an armature, which can slide along the rails. Current is passed through the rails and armature, in the direction of the arrows. There is a magnetic field pointing into the page between the rails, as shown. 

\vspace{2cm}

\textbf{a.} What direction is the force on the armature here?

\textbf{b.} Let's say for simplicity that B is constant and has a magnitude of 2.4 T. Also, the armature length is 50 cm = 0.5 m. Say we want force of magnitude 1.44 million N = $1.44\times 10^6$ N on the armature. How much current do we need? (It's a lot.)

\textbf{c.} To create such a current, the rails and armature are connected to a source of voltage, making a circuit. To allow a larger current for a fixed voltage, should we use thicker rails or thinner rails? Defend your answer. Remember, we need small resistance to get a large current.


\textbf{d.} In fact, B is created by the currents in the rails, so it is not constant. Explain why, using an equation or conceptually.

\textbf{e*.} If we load a more massive projectile into the gun, it will take more energy to accelerate it (say, to a given speed). Yet the force is determined by the current, which is not obviously related to the projectile mass. So how is it that a more massive projectile makes the voltage source provide more power during operation? Hint: think about a motor and back voltage.

\newpage

\textbf{4.} The following diagram shows a coil in a magnetic field.

\vspace{2cm}

\textbf{a.} If we want this coil to act as a generator, generating AC, what do we need to do?

\textbf{b.} If we want the coil to act as a DC motor, what do we need to do?

\textbf{c.} In the generator case, explain specifically why more energy is required to operate the generator if we try to run more electrical appliances with it. 

\bigskip

\textbf{5.} The following is a diagram of a circuit where a voltage source is connected to two resistors in parallel. 

\vspace{2.2cm}

\textbf{a.} If the voltage is 120 V and the resistors are 40 $\Omega$ and 30 $\Omega$, find the current in each resistor.

\textbf{b.} What's the total power being delivered to this circuit by the battery (or whatever the power source is)?

\textbf{c.} Say the electricity is actually AC. If the resistors are, say, a toaster and an electric heater, explain how they can still work even though the voltage and current are constantly changing back and forth. 

\bigskip

\textbf{6a and b.} Choose \textbf{two} of the following \textbf{three} to answer: 
\begin{itemize}
    \item Explain the difference between a paramagnetic material and a ferromagnetic material
    \item Explain how an electromagnet works, including a drawing; 
    \item Explain why resistors like lightbulbs draw more current right when are switched on than after they're been on for a while.
\end{itemize}

\textbf{7.} I put my hand near a doorknob and a spark passes between them, shocking me.

\textbf{a.}  What happened in the air to allow the spark to jump, i.e., charge to flow in it? (It's the same process as when lightning strikes.)

\textbf{b.} As my hand nears the doorknob, the electric field strength between the two increases because opposite charge builds up on each surface. Given that my hand started with some excess negative charge, \textbf{(i)} why does this happen? \textbf{(ii)} what moves in the doorknob - electrons, protons, ions, or a combination? \textbf{(iii)} what moves in my hand - electrons, protons, ions, or a combination?

\textbf{c.} Model my hand and the doorknob as flat surfaces. Give them a reasonable surface area and calculate the minimum amount of charge on each which would allow a spark to jump. 

\end{document}